% Blue River Dam — Research Question
% PhD Thesis Chapter: Lifecycle Authority and Governance Engine

\section{Research Question}
\label{sec:blue-river-dam-question}

\subsection{Problem Statement}
\label{sec:blue-river-dam-question-problem}

Process governance in enterprise systems is typically expressed as policy
documents, configuration files, or runtime checks embedded in application
logic. These mechanisms share a structural defect: the policy and the process
it governs occupy the same mutable runtime space. A process can violate its
declared governance model and continue executing; the violation is detected
only if a monitoring agent happens to observe it. The declaration and the
enforcement are separated by an unbounded observability gap.

This problem becomes acute in autonomic systems where the process itself is
responsible for detecting and remediating its own deviations. If the governance
model that defines what constitutes a deviation can be bypassed or drifted, the
entire autonomic loop loses its normative anchor. The system may emit receipts,
transition states, and report fitness scores while operating in a regime that
no human approver ever authorized.

The specific question this project addresses is: \emph{Can process lifecycle
governance be encoded as compile-time constraints such that illegal state
transitions, authority boundary violations, and unsafe memory operations are
structurally impossible --- not merely detected after the fact?}

\subsection{Old-Regime Assumption Challenged}
\label{sec:blue-river-dam-question-old-regime}

The old-regime assumption is that governance is a runtime concern --- that
process compliance is checked by validators, monitors, and audit logs that
observe an already-running process and report deviations as exceptions. Under
this assumption, a process is presumed compliant until proven otherwise. The
burden of proof falls on the monitoring infrastructure, which can be unavailable,
delayed, or incomplete.

Blue River Dam challenges this assumption at its foundation. The thesis
claim --- supported by the compile-time evidence documented in
\texttt{GENERATION\_RECEIPT.md} --- is that governance rules can be
encoded as Rust type constraints such that a binary that successfully compiles
cannot, by construction, execute an illegal transition. Compliance is not
checked; it is \emph{required by the language}. The old-regime monitoring
infrastructure becomes a secondary audit layer over an already-constrained
execution environment, rather than the primary enforcement mechanism.

\subsection{Hypothesis}
\label{sec:blue-river-dam-question-hypothesis}

A Rust library manufactured from formal governance doctrine, using typestate
encoding and \texttt{\#![forbid(unsafe\_code)]}, will satisfy three
properties simultaneously: (1) every lifecycle state transition is guarded
by a named quality gate whose criterion is expressed in the source as a
typed predicate; (2) every MAPE-K component produces a typed artifact
(Evidence, Analysis, Plan, Receipt, Knowledge) that downstream components
consume by type rather than by convention; and (3) the compiled binary is
incapable of executing a path that would leave the system in an undocumented
state. These three properties together constitute compile-time governance
closure.

\subsection{Success Criteria}
\label{sec:blue-river-dam-question-success}

Success is measured against four criteria, all of which are verified by the
artifacts in this project:

\begin{enumerate}
  \item \textbf{Zero unsafe code.} The \texttt{\#![forbid(unsafe\_code)]}
    attribute is present in \texttt{src/lib.rs} and compilation succeeds,
    as certified by \texttt{GENERATION\_RECEIPT.md}.
  \item \textbf{MAPE-K loop closure.} All five loop-closure conditions
    (typed Evidence, typed Analysis, typed Plan, typed Receipt, replayable
    Knowledge) are satisfied, as documented in the receipt.
  \item \textbf{All six quality gates implemented.} Gates covering Design,
    Simulation, Monitoring, Repair, Optimization, and Decommissioning states
    are callable and gate-guarded in \texttt{transition\_state()}.
  \item \textbf{Tests pass.} The initial 5/5 receipt-certified tests pass;
    the 3 post-receipt tests (\texttt{test\_validate\_wf\_net\_soundness},
    \texttt{test\_compute\_fitness\_genuine},
    \texttt{test\_mape\_k\_cycle\_execution}) exercise the deeper
    WF-net soundness and fitness computation paths whose pass/fail status
    is not yet captured in a receipt (see
    Section~\ref{sec:blue-river-dam-questions-gaps}).
\end{enumerate}
